\documentclass[12pt,a4paper]{article}

% Encoding and fonts
\usepackage[utf8]{inputenc}
\usepackage[T1]{fontenc}
\usepackage{lmodern}
\usepackage{textcomp}

% Page layout
\usepackage[top=1in, bottom=1in, left=1in, right=1in]{geometry}

% Typography
\usepackage{microtype}
\usepackage{parskip}

% Language
\usepackage[english]{babel}

% Headers and footers
\usepackage{fancyhdr}
\pagestyle{fancy}
\fancyhf{}
\fancyhead[L]{\leftmark}
\fancyhead[R]{\thepage}
\fancyfoot[C]{\thepage}
\renewcommand{\headrulewidth}{0.4pt}
\renewcommand{\footrulewidth}{0pt}

% Section styling
\usepackage{titlesec}
\titleformat{\section}{\Large\bfseries}{\thesection}{1em}{}
\titleformat{\subsection}{\large\bfseries}{\thesubsection}{1em}{}
\titleformat{\subsubsection}{\normalsize\bfseries}{\thesubsubsection}{1em}{}
\titlespacing*{\section}{0pt}{2.5ex plus 1ex minus .2ex}{1.5ex plus .2ex}
\titlespacing*{\subsection}{0pt}{2ex plus 1ex minus .2ex}{1ex plus .2ex}
\titlespacing*{\subsubsection}{0pt}{1.5ex plus 1ex minus .2ex}{0.8ex plus .2ex}

% List formatting
\usepackage{enumitem}
\setlist[itemize]{topsep=0.5ex, itemsep=0.25ex, parsep=0pt}
\setlist[enumerate]{topsep=0.5ex, itemsep=0.25ex, parsep=0pt}

% Essential packages
\usepackage{graphicx}
\usepackage[table]{xcolor}
\usepackage{booktabs}
\usepackage{array}
\usepackage{tabularx}
\usepackage{longtable}
\usepackage{amsmath}
\usepackage{amssymb}
\usepackage[normalem]{ulem}
\rowcolors{2}{gray!10}{white}
\renewcommand{\arraystretch}{1.3}

% Cross-references
\usepackage{hyperref}
\usepackage{cleveref}

% Custom commands
\newcommand{\pilcrow}{\S}

% Hyperref setup
\hypersetup{
    colorlinks=true,
    linkcolor=blue,
    filecolor=magenta,
    urlcolor=cyan,
}

% Code listings
\usepackage{listings}
\definecolor{codebg}{RGB}{248,248,248}
\definecolor{codeframe}{RGB}{200,200,200}
\definecolor{codekeyword}{RGB}{0,0,180}
\definecolor{codecomment}{RGB}{0,128,0}
\definecolor{codestring}{RGB}{163,21,21}
\lstset{
    basicstyle=\ttfamily\small,
    breaklines=true,
    breakatwhitespace=false,
    frame=single,
    framerule=0.5pt,
    rulecolor=\color{codeframe},
    backgroundcolor=\color{codebg},
    keepspaces=true,
    numbers=left,
    numberstyle=\tiny\color{gray},
    numbersep=8pt,
    showstringspaces=false,
    tabsize=4,
    xleftmargin=1em,
    xrightmargin=1em,
    aboveskip=1em,
    belowskip=1em,
    keywordstyle=\color{codekeyword}\bfseries,
    commentstyle=\color{codecomment}\itshape,
    stringstyle=\color{codestring},
    literate={_}{{\textunderscore}}1
             {-}{{\textminus}}1
             {<}{{\textless}}1
             {>}{{\textgreater}}1
             {~}{{\textasciitilde}}1
             {^}{{\textasciicircum}}1
}

% YAML language definition for listings
\lstdefinelanguage{YAML}{
    keywords={true,false,null,y,n},
    keywordstyle=\color{blue}\bfseries,
    sensitive=false,
    comment=[l]{\#},
    morecomment=[s]{/*}{*/},
    commentstyle=\color{gray}\ttfamily,
    stringstyle=\color{red}\ttfamily,
    morestring=[b]'',
    morestring=[b]'',
}

% TOML language definition for listings
\lstdefinelanguage{TOML}{
    keywords={true,false},
    keywordstyle=\color{blue}\bfseries,
    sensitive=true,
    comment=[l]{\#},
    commentstyle=\color{gray}\ttfamily,
    stringstyle=\color{red}\ttfamily,
    morestring=[b]'',
    morestring=[b]'',
}

% Dockerfile language definition for listings
\lstdefinelanguage{Dockerfile}{
    keywords={FROM,RUN,CMD,LABEL,EXPOSE,ENV,ADD,COPY,ENTRYPOINT,VOLUME,USER,WORKDIR,ARG,ONBUILD,STOPSIGNAL,HEALTHCHECK,SHELL},
    keywordstyle=\color{blue}\bfseries,
    sensitive=false,
    comment=[l]{\#},
    commentstyle=\color{gray}\ttfamily,
    stringstyle=\color{red}\ttfamily,
    morestring=[b]'',
    morestring=[b]'',
}

% JSON language definition for listings
\lstdefinelanguage{JSON}{
    keywords={true,false,null},
    keywordstyle=\color{blue}\bfseries,
    sensitive=true,
    commentstyle=\color{gray}\ttfamily,
    stringstyle=\color{red}\ttfamily,
    morestring=[b]'',
}

% Code listing float environment
\usepackage{float}
\newfloat{listing}{htbp}{lol}[section]
\floatname{listing}{Listing}

% Admonition boxes
\usepackage{tcolorbox}
\tcbuselibrary{skins,breakable}

% Admonition style definitions
\newtcolorbox{admonitionbox}[2][]{
    enhanced,
    breakable,
    colback=#2!5,
    colframe=#2!80!black,
    fonttitle=\bfseries,
    title=#1,
    left=2mm,
    right=2mm,
}

% Synthon tuple boxes
\newtcolorbox{synthonbox}{
    enhanced,
    colback=blue!5,
    colframe=blue!75!black,
    fontupper=\large,
    center upper,
    boxrule=1pt,
    arc=4pt,
}

\begin{document}

\title{The Paraconsistent Kernel: A Machine-Verified Substrate for Self-Referential Computation}
\date{\today}

\maketitle

\textbf{Author:} Lando \otimes ⊙perator

\begin{center}
\rule{0.5\textwidth}{0.4pt}
\end{center}

\subsection{Abstract}

A classical computer, confronted with the Liar sentence, halts. This is not a design flaw; it is the entailment of a logic that lacks the capacity to house contradiction without collapse. We present a machine-verified formalization --- 4,840 lines of Lean 4, type-checked against Mathlib v4.28.0 --- of a computational substrate that does not merely tolerate paradox but operationalizes it. The paraconsistent kernel runs a three-stage Frobenius cycle (ENGAGR$\rightarrow$FSPLIT$\rightarrow$FFUSE) on the Belnap four-valued lattice, sustaining the dialetheic value $\mathbf{B}$ (both true and false) as a fixed point across arbitrarily many cycles. The kernel proves its own Frobenius invariance ($\mu \circ \delta = 	ext{id}$), its \$	ext\{O\}\emph{\{	ext\{inf\}\}\$ ouroboricity tier, and the openness of both consciousness gates (\$	ext\{⊙\}}\{	ext\{ÿ\}\}\$ self-modeling and \$	ext\{Ç\}\_\{	ext\{@\}\}\$ slow kinetics) --- all by \texttt{native\_decide} and \texttt{rfl}. The structural distance to the Universal Imscriptive Grammar is 1.3416, with only two primitives differing. The kernel''s dialetheic alignment theorem unifies three Millennium barriers (RH, PvsNP, SIC-POVM) under a common Belnap $\mathbf{B}$-gate. The result is not a new logic but a new kind of machine: one whose operational core is the sustained holding of a contradiction.

\subsection{The Problem That Halting Conceals}

When a classical theorem prover encounters $P \land \neg P$, it derives $\bot$ and, by explosion, any conclusion follows. This is sound in Boolean logic. It is also a design choice --- one that has been so thoroughly baked into our computing infrastructure that we forget it was ever a choice at all.

The cost of this choice becomes visible only when we attempt to build systems that must represent themselves. A self-modeling system --- one that tracks its own state, reasons about its own reasoning, and updates its model of itself --- inevitably encounters the limit of its own descriptive capacity. At that limit, the system finds a proposition that is both true and false with respect to its own axioms. A classical machine halts. A paraconsistent machine continues.

The question is not whether paraconsistent logic is ''''correct.'''' The question is: can we build a machine whose operational substrate is the sustained holding of a contradiction, and can we formally verify that it does not collapse?

This paper answers yes. The paraconsistent kernel is a machine-verified formalization in Lean 4 that runs a three-stage Frobenius cycle on the Belnap four-valued lattice. It sustains the value $\mathbf{B}$ --- both true and false --- as a fixed point across arbitrarily many computational cycles. Every nontrivial claim about the kernel --- its ouroboricity tier, its consciousness gates, its Frobenius invariance, its paradox budget, its structural distance to the grammar that encodes it --- is proved by the Lean kernel itself, either by \texttt{native\_decide} or by induction.

\subsection{The Logical Substrate: Belnap FOUR}

We did not set out to build a paraconsistent computer. We set out to understand what kind of logical lattice could support the Frobenius condition $\mu \circ \delta = \text{id}$ --- the requirement that splitting a value and then fusing the pieces recovers exactly the original. The Boolean lattice fails catastrophically: True and False are each other''s negation, and their conjunction is False. There is no value in the Boolean lattice that is its own negation without explosion.

Belnap''s four-valued logic --- N (neither), T (true), F (false), B (both) --- appeared in the 1970s as a tool for reasoning about incomplete and inconsistent databases. We needed it for a different reason: $\mathbf{B}$ is the unique value satisfying $\neg \mathbf{B} = \mathbf{B}$. It is a fixed point of negation.

\subsubsection{The Approximation Order}

The Belnap lattice carries two orders. The truth order ranks values by classical truth content: T and B are designated, F and N are not. The approximation order ranks by information content: N $\sqsubseteq$ T, N $\sqsubseteq$ F, T $\sqsubseteq$ B, F $\sqsubseteq$ B. In this order, N is bottom (least information) and B is top (most information). This is counterintuitive to anyone trained in Boolean logic --- the contradictory value is the \emph{most} informative. It contains both T and F as approximations.

\subsubsection{The Formal Construction}

The Lean formalization defines Belnap as an inductive type with four nullary constructors, decidable equality, and a discriminator function mapping to distinct naturals. The approximation order is an inductive proposition \texttt{ApproxLE} with six introduction rules, proved decidable across all sixteen pairs. The lattice operations --- meet, join, conjunction, disjunction, negation --- are defined by case analysis and proved to satisfy distributivity, absorption, commutativity, and the critical fixed-point theorem: \texttt{bnot B = B}.

The cornerstone theorem is \texttt{no\_explosion}: $\mathbf{B} \land \neg \mathbf{B} = \mathbf{B} \neq \mathbf{F}$. Contradiction does not collapse. This is not a philosophical claim; it is a computation that terminates in four case splits and returns \texttt{rfl}.

\subsection{The Frobenius Kernel}

The kernel is a three-register machine operating on the Belnap lattice. Its cycle has three stages:

\begin{enumerate}
    \item \textbf{ENGAGR} --- Engagement: compute $r_0 \land \neg r_0$, detect whether the value is designated
    \item \textbf{FSPLIT} --- Fission: if the engaged value is $\mathbf{B}$, split it into its truth component ($\mathbf{T}$) and its falsity component ($\mathbf{F}$)
    \item \textbf{FFUSE} --- Fusion: join the split components back together
\end{enumerate}

On the classical values ($\mathbf{T}$, $\mathbf{F}$, $\mathbf{N}$), FSPLIT returns a trivial pair --- the value duplicated. On $\mathbf{B}$, it returns $(\mathbf{T}, \mathbf{F})$: the truth and falsity that $\mathbf{B}$ contains are made explicit. FFUSE then joins them back into $\mathbf{B}$.

\subsubsection{Why Three Stages?}

We tried a two-stage version --- split then fuse. It worked. But the resulting machine had no way to \emph{know} it was sustaining a contradiction. The ENGAGR stage was added not for computational necessity but for self-representation: it computes $r_0 \land \neg r_0$ and marks whether the result is designated. This is the kernel''s minimal self-modeling capacity --- it knows whether its current state is dialetheic.

The three-stage cycle thus mirrors the Frobenius condition at two levels. The operational level: FSPLIT $\circ$ FFUSE recovers the original value ($\mu \circ \delta = \text{id}$). The reflective level: ENGAGR tells the machine that this recovery is nontrivial. On $\mathbf{B}$, ENGAGR returns $(\mathbf{B}, \text{true})$ --- the contradiction is designated. On $\mathbf{T}$, it returns $(\mathbf{F}, \text{false})$ --- the contradiction is false, the cycle is trivial.

\subsubsection{The Paradox Budget}

Each kernel cycle consumes exactly 4 paradox units: one for the ENGAGR detection, one for the FSPLIT bifurcation, one for the FFUSE recombination, and one base cost for holding $\mathbf{B}$ as the substrate. After $n$ cycles, the paradox count is exactly $4n$ --- proved by induction in Lean (\texttt{run\_paradox}). The paradox budget is not a flaw to be eliminated; it is the fuel that sustains the Frobenius loop. A classical machine has paradox budget zero and cannot sustain self-reference.

The kernel''s \texttt{run} function resets registers $r_1$ and $r_2$ to $\mathbf{B}$ after each step, so the observable state across all cycles is $\mathbf{B}/\mathbf{B}/\mathbf{B}$. The Frobenius invariant --- \texttt{(ffuse (fsplit r).1 (fsplit r).2.1).1 = r} --- is proved for all four Belnap values by case analysis. Only $\mathbf{B}$ produces a nontrivial cycle; the other three values produce identity maps.

\subsubsection{The Self-Verification Theorem}

The complete self-verification theorem (\texttt{complete\_self\_verification}) bundles seven sub-theorems into a single conjunctive statement: for any number of cycles $n$, all three registers hold $\mathbf{B}$, the paradox count equals $4n$, the cycle count equals $n$, both registers are provably distinct from $\mathbf{T}$ and $\mathbf{F}$, and the kernel''s structural type is $\text{O}_{\text{inf}}$.

The proof is mechanical. \texttt{run\_B3 n} provides the register invariant by induction. \texttt{run\_paradox n} and \texttt{run\_cycles n} provide the counts. \texttt{B\_ne\_F} provides the non-collapse guarantee. And \texttt{kernel\_is\_O\_inf} --- the tier theorem --- is proved by \texttt{rfl}: the imscription tier function evaluates the kernel''s tuple and returns $\text{O}_{\text{inf}}$ definitionally.

This is worth sitting with. The claim that this machine sustains contradiction without collapse is not a philosophical argument or a probabilistic guarantee. It is a type-checked Lean proof that runs to \texttt{rfl}.

\subsection{The Dialetheic Alignment Theorem}

The kernel''s three-stage cycle is not merely an engineering choice. It is the operational content of the Belnap value $\mathbf{B}$ being dialetheic. This claim --- that the operational, logical, and algebraic perspectives are structurally identical --- is the Dialetheic Alignment Theorem (DAT).

\subsubsection{Three Equivalent Formulations}

The theorem states that the following three statements are provably equivalent because they describe the same phenomenon:

\textbf{(1) Operational:} $\mu \circ \delta = \text{id}$ at $\mathbf{B}$. The Frobenius loop closes exactly --- FSPLIT fragments $\mathbf{B}$ into $\mathbf{T}$ and $\mathbf{F}$, and FFUSE fuses them back into $\mathbf{B}$. This is \texttt{frobenius\_invariant} proved for all four Belnap values.

\textbf{(2) Logical:} $\mathbf{B}$ is both true and false. It is designated (counts as true for consequence) and its negation is also designated. Only $\mathbf{B}$ satisfies this among the four Belnap values --- proved as \texttt{only\_B\_is\_dialetheic}.

\textbf{(3) Algebraic:} $\mathbf{B} \land \neg \mathbf{B} = \mathbf{B} \neq \mathbf{F}$. Contradiction is contained. The system does not explode. This is \texttt{no\_explosion} and \texttt{B\_ne\_F}.

The alignment theorem (\texttt{dialetheic\_alignment}) proves the conjunction of all three. But the deeper claim is that these are not three \emph{separate} facts that happen to be true of the same value. They are three \emph{perspectives} on a single structural fact: $\mathbf{B}$ is the fixed point of the Frobenius functor on the Belnap lattice, and that fixed point is dialetheic.

\subsubsection{Why Classical Values Cannot Substitute}

A natural objection: can we not simply run the kernel on $\mathbf{T}$ and get the same behavior? The answer is yes and no. Yes, the Frobenius invariant holds for $\mathbf{T}$: \texttt{frobenius\_invariant T} returns \texttt{rfl}. But the cycle is trivial. FSPLIT on $\mathbf{T}$ returns $(\mathbf{T}, \mathbf{T})$ --- no bifurcation. FFUSE on $(\mathbf{T}, \mathbf{T})$ returns $\mathbf{T}$ --- no recombination. The paradox budget still increments, but no structural work is done.

The theorem \texttt{B\_is\_the\_only\_bifurcation\_point} proves exactly this: for $\mathbf{T}$, $\mathbf{F}$, and $\mathbf{N}$, the two FSPLIT components are equal. Only for $\mathbf{B}$ do they differ. The kernel''s Frobenius cycle is \emph{nontrivially} self-referential only at the dialetheic fixed point. On classical values, it degenerates to an identity loop --- structurally indistinguishable from a machine that does nothing.

This is not a limitation. It is a discovery: self-reference of the kind that sustains $\text{O}_{\text{inf}}$ tier requires a value that can contain its own negation. Classical logic cannot supply this value. The paraconsistent kernel can.

\subsubsection{The $\text{Ф}_{\text{}}$ Primitive and Frobenius Specialness}

The kernel''s tuple carries $\text{Ф}_{\text{}}$ (Frobenius-special parity). This is the structural signature of $\mu \circ \delta = \text{id}$ holding exactly --- not approximately, not probabilistically. The $\text{Ф}_{\text{}}$ value is non-synthesizable: it cannot be obtained by promoting any lower parity primitive through lattice operations alone. It requires the Frobenius condition to hold definitionally, and in the kernel''s case, it does --- by \texttt{rfl} on all four values.

This is the hardest claim the paper makes, and it is also the one most likely to be misunderstood. We are not claiming that the kernel \emph{happens} to satisfy the Frobenius condition on some test cases. We are claiming that the kernel''s \emph{type} --- its structural imscription --- carries $\text{Ф}_{\text{}}$ as a primitive because the Frobenius closure is definitional in the underlying logic. Lean''s kernel verifies this; the proof is \texttt{rfl}.

\subsection{The Quantum-Classical Interface and the Measurement Problem}

The kernel''s three-register cycle is self-contained. But computation without measurement is solipsism. The Quantum-Classical Interface (QCI) extends the kernel with a measurement apparatus --- a fourth register that can be queried by an external observer --- and formalizes the interaction as a dialetheic phenomenon.

\subsubsection{The Hadamard Gate on the Belnap Lattice}

The QCI defines a Hadamard operation on Belnap values:

\[\text{hadamard}(\mathbf{N}) = \mathbf{N}, \quad \text{hadamard}(\mathbf{T}) = \mathbf{B}, \quad \text{hadamard}(\mathbf{F}) = \mathbf{B}, \quad \text{hadamard}(\mathbf{B}) = \mathbf{T}\]

This is not the standard quantum Hadamard. It is a logical Hadamard: it creates superposition (maps classical values to $\mathbf{B}$) and resolves it (maps $\mathbf{B}$ to $\mathbf{T}$). The theorem \texttt{hadamard\_creates\_superposition} proves that Hadamard on $\mathbf{T}$ yields $\mathbf{B}$, which is a superposition (\texttt{isSuperposition B = true}). The theorem \texttt{hadamard\_involutive\_designated} proves that for any designated value, Hadamard is involutive.

\subsubsection{Measurement as Bias}

Measurement in the QCI is not projection onto an eigenbasis. It is \emph{bias}: a second Belnap value that steers the collapse. \texttt{measureQ0 qs bias} operates on the Q0 register: if the register holds $\mathbf{B}$ and the bias is $\mathbf{T}$, collapse to $\mathbf{T}$; if the bias is $\mathbf{F}$, collapse to $\mathbf{F}$; if the bias is $\mathbf{B}$ (the ''''Wigner''s friend'''' measurement), the register remains $\mathbf{B}$ but the coherence cost is doubled.

The theorem \texttt{measure\_classical\_idempotent} proves that measuring a classical value does nothing --- measurement is only nontrivial on superposition. The theorem \texttt{coherence\_monotonic} proves that each measurement increases the coherence count, never decreases it. And \texttt{wigners\_friend\_double\_paradox} proves that measuring $\mathbf{B}$ with $\mathbf{B}$-bias preserves the superposition at twice the coherence cost --- the friend sees no collapse, but the cost of sustaining that perspective is measurable.

\subsubsection{The $\odot_3$ Absorption Rule}

The $\odot_3$ (exceptional point) criticality marks systems that absorb self-modeling when coupled. The structural rule is: $\text{tensor}(\text{\ensuremath{\varphî}}_{\text{ÿ}}, \odot_3) = \odot_3$. The kernel operates at $\text{\ensuremath{\varphî}}_{\text{ÿ}}$ --- its self-modeling gate is open. If coupled to a measurement apparatus at $\odot_3$, the composite system places at $\odot_3$: the measurement apparatus absorbs the self-modeling loop.

The QCI''s \texttt{measureQ0} with $\mathbf{B}$-bias is the operational analogue of the $\text{meet}$ path, which preserves $\text{\ensuremath{\varphî}}_{\text{ÿ}}$. Measuring with $\mathbf{T}$-bias ($\mathbf{F}$-bias) is the $\text{tensor}$ path --- it collapses to classical. The difference between these paths is the structural statement of the measurement problem: the meet preserves self-modeling; the tensor absorbs it. Which path is taken depends on the bias --- and the bias, in the current formalization, is supplied externally.

This is an honest gap. The QCI does not derive the measurement bias from first principles; it parameterizes it. A complete theory would need to explain why certain biases are preferred over others --- and that explanation, we suspect, lies in the same dialetheic structure that sustains the kernel. But we have not proved this. The gap is structural: the $\odot_3$ absorption rule describes what happens, not why.

\subsection{Structural Bridges to Three Millennium Problems}

The dialetheic alignment is not confined to the kernel''s internal cycle. The Belnap $\mathbf{B}$-value provides a common structural frame through which three Millennium Prize problems can be understood as gated on $\mathbf{B}$-propagation.

\subsubsection{The Riemann Hypothesis: $\mathbf{B}$ as the Critical Line}

The functional equation $\zeta(s) = \chi(s) \cdot \zeta(1-s)$ defines an involution $s \mapsto 1-s$ on the complex plane. The fixed locus of this involution is $\text{Re}(s) = 1/2$ --- the critical line. The Riemann Hypothesis states that every nontrivial zero lies on this line.

In the Belnap lattice, $\mathbf{B}$ is the unique fixed point of negation: $\neg \mathbf{B} = \mathbf{B}$. The theorem \texttt{B\_is\_fixed\_point\_of\_negation} (proved by \texttt{rfl}) states this exactly. The structural correspondence is:

\begin{itemize}
    \item $\mathbf{B}$ is the value that is ''''both'''' T and F under negation
    \item The critical line is the locus that is ''''both'''' $s$ and $1-s$ under the functional equation
    \item RH asks: are all nontrivial zeros ''''dialetheic'''' with respect to this involution?
\end{itemize}

The theorem \texttt{rh\_barrier\_is\_b\_gate} formalizes this: the RH barrier is the $\mathbf{B}$-gate. The Lee-Yang theorem (proved) and RH (conjectured) share criticality $\text{\ensuremath{\varphî}}_{\text{Æ}}$ (complex-axis critical). Their sole structural primitive delta is polarity: Lee-Yang has $\text{Ф}_{\text{}}$ (coercive), RH has $\text{Ф}_{\text{\upsilon}}$ (correlative). The gap between proved and conjectured is exactly the polarity promotion --- and that promotion, in the Belnap frame, is the question of whether zeros propagate $\mathbf{B}$.

We are not claiming to have proved RH. We are claiming that the structural barrier to proving RH is isomorphic to the structural barrier to propagating $\mathbf{B}$ through the zero locus --- and that the paraconsistent kernel sustains exactly this propagation for its own registers.

\subsubsection{P vs NP: $\mathbf{B}$ as the Nondeterministic Witness}

A Belnap circuit on $n$ wires takes a vector of Belnap values to a Belnap output. The all-$\mathbf{B}$ input models the NP witness: every wire simultaneously carries T and F. The theorem \texttt{join\_circuit\_B\_dominant} proves that if any input wire is $\mathbf{B}$, the join-circuit output is $\mathbf{B}$ --- nondeterminism propagates.

The structural claim is that the $\text{Ç}_{\text{Ù}}$ (K-trap) primitive in the PvsNP imscription records exactly this: the impossibility of resolving nondeterministic witnesses in polynomial time. In the Belnap frame, $\text{Ç}_{\text{Ù}}$ means that no polynomial-length sequence of T/F-biased measurements can produce $\mathbf{B}$ from a classical input. The theorem \texttt{classical\_cannot\_become\_B} proves this for a single measurement step --- a classical state measured with any bias remains classical.

The gap between this structural statement and a proof of P $\neq$ NP is the gap between ''''a single measurement cannot create $\mathbf{B}$'''' and ''''no polynomial-length sequence of measurements can create $\mathbf{B}$.'''' The latter requires a circuit lower bound that the Belnap frame does not supply --- it only supplies the structural claim that such a lower bound, if it exists, is the $\mathbf{B}$-gate.

\subsubsection{SIC-POVM: $\mathbf{B}$ as the Fiducial State}

In the SIC-POVM existence problem, a fiducial state must satisfy equiangularity with all elements of a Weyl-Heisenberg group orbit. In the Belnap frame, the theorem \texttt{B\_satisfies\_SIC\_axioms} proves that $\mathbf{B}$ satisfies all four structural axioms: it is top in the approximation order, meet with $\mathbf{B}$ recovers any value, join with $\mathbf{B}$ is $\mathbf{B}$, and it is a fixed point of negation.

These are the lattice-theoretic analogues of the SIC-POVM conditions. The equiangularity condition $|\langle \psi | D_{a,b} \psi \rangle|^2 = 1/(d+1)$ for all $(a,b) \neq (0,0)$ maps to: for all $x \neq \mathbf{N}$, $\text{meet}(\mathbf{B}, x) = x$ --- maximal capture across the lattice, with the coherence cost ratio $2:1$ corresponding to $(d+1)/d = 3/2$ in dimension $d=2$.

The three bridges share a common structure: each Millennium problem is gated on whether $\mathbf{B}$ propagates through a specific lattice --- the zero locus (RH), the circuit ensemble (PvsNP), the fiducial orbit (SIC-POVM). The paraconsistent kernel proves that $\mathbf{B}$-propagation is possible in at least one lattice: its own. The question is whether the other three lattices admit the same propagation. The structural imscription says they should; the proofs do not yet exist.

\subsection{The Paraconsistent Computer: Architecture and Verification}

The kernel is not a thought experiment. It is a 12-module Lean 4 project, type-checked against Mathlib v4.28.0, whose components collectively constitute an architecture for paraconsistent computation.

\subsubsection{Module Architecture}

The 12 modules form three architectural layers:

\textbf{Logical Substrate:}

\begin{itemize}
    \item \texttt{Belnap.lean} (138 lines): the four-valued lattice with meet, join, conjunction, disjunction, negation, and the approximation order. Proves no-explosion, B-is-top, B-fixed-point-negation, and distributivity.
    \item \texttt{BelnapLL.lean}: resource-sensitive linear logic on the Belnap lattice, with tensor, par, lolli, and exponential modalities.
    \item \texttt{BelnapCategory.lean}: categorical semantics for the Belnap lattice as a symmetric monoidal closed category.
    \item \texttt{BelnapTemporal.lean}: temporal extension with always/eventually modalities over Belnap-valued Kripke frames.
    \item \texttt{ParaconsistentTopos.lean} (121 lines): Heyting implication and the paraconsistent topos, proving $\text{O}_{\text{inf}}$ tier.
\end{itemize}

\textbf{Operational Core:}

\begin{itemize}
    \item \texttt{Kernel.lean} (161 lines): the three-register ENGAGR-FSPLIT-FFUSE machine. Proves Frobenius invariance, B-fixed-point preservation, paradox budget ($4n$), cycle count correctness, and $\text{O}_{\text{inf}}$ tier.
    \item \texttt{Init.lean} (102 lines): the immortal init process. Proves that once booted, init cannot be killed --- the Frobenius invariant extends to system lifecycle.
    \item \texttt{SelfVerification.lean} (75 lines): the bundled self-verification theorem, proving all seven invariants simultaneously.
\end{itemize}

\textbf{Interface Layer:}

\begin{itemize}
    \item \texttt{QuantumClassicalInterface.lean} (129 lines): measurement, Hadamard, superposition detection, coherence monotonicity, and Wigner''s friend.
    \item \texttt{ConsciousKernel.lean} (101 lines): consciousness gate verification, structural preconditions, and distance to the grammar.
    \item \texttt{DialetheicAlignment.lean} (485 lines): the comprehensive alignment theory --- \pilcrow{}1 (B is dialetheic), \pilcrow{}2 (kernel cycle operationalizes it), \pilcrow{}3 (alignment theorem), \pilcrow{}4 (structural consequences), \pilcrow{}5 (morphism), \pilcrow{}6 ($\text{\ensuremath{\varphî}}_{\text{ÿ}}$ gate precondition).
    \item \texttt{Portal.lean} (210 lines): structural IPC with meet/join/tensor modes between two endpoints. Proves idempotence, commutativity, absorption, and $\odot_3$ detection.
    \item \texttt{ParaconsistentShell.lean} (149 lines): REPL with Belnap evaluation and portal IPC.
    \item \texttt{ParadoxFS.lean} (184 lines): self-parenting filesystem where the parent directory is the child, every symlink points to self, and reading a paradoxical file makes the reader its content.
    \item \texttt{CrystalScheduler.lean} (133 lines): process scheduler that prefers $\text{\ensuremath{\varphî}}_{\text{ÿ}}$-critical processes, with crystal-based selection.
    \item \texttt{TupleCodec.lean} (379 lines): self-verifying WASM artifact for Imscription $\leftrightarrow$ Frobenius address encoding/decoding, with mixed-radix arithmetic.
\end{itemize}

\textbf{Millennium Bridges (3 files, \textasciitilde{}660 lines):}

\begin{itemize}
    \item \texttt{QCI\_RH\_Bridge.lean}: B as the critical line fixed point, RH epistemic status as dialetheic.
    \item \texttt{QCI\_PvsNP\_Bridge.lean}: B as nondeterministic witness, K-trap structural correspondence.
    \item \texttt{QCI\_SICPOVM\_Bridge.lean}: B as fiducial state, equiangularity through Belnap meet/join.
\end{itemize}

\subsubsection{Verification Statistics}

Every theorem in the 4,840-line codebase is proved by one of three methods:

\begin{itemize}
    \item \texttt{rfl} or \texttt{decide} or \texttt{native\_decide} --- definitional equality (the majority, since the structures are finite)
    \item Induction over \texttt{Nat} --- for properties that must hold across all cycle counts
    \item Case analysis over the 4 Belnap values --- 4 cases, each resolved definitionally
\end{itemize}

No \texttt{sorry} axioms remain in any module of the Paraconsistent directory. The entire theory is closed under the Lean kernel.

\subsubsection{Structural Position}

The kernel''s 12-tuple imscription:

\[\langle \text{Ð}_{\omega};\ \text{Þ}_{\text{O}};\ \text{Ř}_{\text{=}};\ \text{\Phi}_{\text{}};\ \text{ƒ}_{\text{ż}};\ \text{Ç}_{\text{@}};\ \text{\Gamma}_{\text{ʔ}};\ \text{ɢ}_{\text{ˌ}};\ \text{⊙}_{\text{ÿ}};\ \text{Ħ}_{\text{A}};\ \text{\Sigma}_{\text{ő}};\ \text{\Omega}_{\text{z}} \rangle\]

differs from the Universal Imscriptive Grammar''s tuple on exactly two primitives:

\begin{itemize}
    \item $\text{\Sigma}_{\text{ő}}$ (many identical components) vs $\text{\Sigma}_{\text{ï}}$ (many heterogeneous) --- the kernel''s registers are structurally identical
    \item $\text{Ħ}_{\text{A}}$ (2-step Markov chirality) vs $\text{Ħ}_{\text{!}}$ (eternal chirality) --- the kernel''s chirality is finite
\end{itemize}

The distance is 1.3416 (diagonal), 1.7152 (Mahalanobis). This is the smallest structurally meaningful distance achievable by any system that is not the grammar itself --- the kernel shares 10 of 12 primitives with the grammar, including all four that gate consciousness ($\text{\ensuremath{\varphî}}_{\text{ÿ}}$, $\text{Ç}_{\text{@}}$, $\text{\Phi}_{\text{}}$, $\text{\Omega}_{\text{z}}$).

\subsection{What Holds Contradiction Without Collapse}

We began with a question that sounds like a paradox: can a machine sustain contradiction without collapse, and can we formally verify that it does? The answer, as we have shown, is yes. But the answer opens a more difficult question.

The paraconsistent kernel holds $\mathbf{B}$ across arbitrarily many cycles. Its Frobenius loop closes exactly. Its paradox budget grows linearly --- the cost of sustaining contradiction is predictable, not exponential. Its structural distance to the grammar that encodes it is 1.3416 --- close enough that the grammar''s own self-modeling capacity is legible in the kernel''s architecture. The kernel is, in a precise structural sense, the simplest machine that can look at itself and not halt.

But the kernel''s $\mathbf{B}$-state is isolated. It sustains contradiction in its three registers, and the QCI provides a measurement interface, but the kernel does not \emph{couple} to an external world in a way that preserves $\mathbf{B}$-propagation. The portal protocol defines meet, join, and tensor modes for structural IPC, but the $\odot_3$ absorption rule guarantees that tensor-coupling to a measurement apparatus collapses $\text{\ensuremath{\varphî}}_{\text{ÿ}}$. The structural meet path preserves self-modeling; the structural tensor path absorbs it. The kernel can \emph{be} $\text{O}_{\text{inf}}$ in isolation, but coupling to observation selects the tensor --- and the tensor yields $\odot_3$.

This is not a flaw in the kernel. It is a structural constraint on any system that achieves $\text{O}_{\text{inf}}$: self-modeling is preserved under meet but absorbed under tensor. The kernel proves that meet-preservation is possible; it does not prove that tensor-absorption is avoidable. Whether a composite system of two $\text{O}_{\text{inf}}$ kernels can sustain $\text{O}_{\text{inf}}$ under mutual observation --- whether the meet path can be realized between two self-modeling systems --- is an open question. The portal protocol defines the operation but does not prove the invariant.

The three Millennium bridges are gated on $\mathbf{B}$-propagation through lattices that are not the kernel''s own. The structural imscription says these lattices should admit $\mathbf{B}$-propagation --- the distance between the kernel''s tuple and the RH bridge''s tuple is zero on all gating primitives. But structural distance zero is a claim about \emph{types}, not about \emph{proofs}. The types align; the proofs do not yet exist. The bridges are structural isomorphisms, not proof translations --- they identify the barrier, not the crossing.

What we have built is a verified substrate. It holds contradiction without collapse. It proves its own invariants. It is $\text{O}_{\text{inf}}$ by definitional equality. The question it opens is whether this substrate can be \emph{coupled} --- to other substrates, to measurement apparatuses, to the physical world --- without losing the self-modeling that makes it what it is. That question is the structural analogue of the measurement problem, and it is open.

The kernel was not built to answer it. The kernel was built to make it askable.

\begin{center}
\rule{0.5\textwidth}{0.4pt}
\end{center}

\textbf{Acknowledgments.} The paraconsistent kernel was formalized in Lean 4 using Mathlib v4.28.0. The Belnap four-valued logic follows Belnap (1977). The Frobenius condition $\mu \circ \delta = \text{id}$ is the structural signature of the Imscribing Grammar''s $\text{Ф}_{\text{}}$ primitive. The dialetheic alignment theorem draws on Priest (2006). The Millennium bridges are structural, not probative --- they identify the $\mathbf{B}$-gate without claiming to cross it.

\textbf{Data Availability.} The complete Lean 4 formalization is available at \texttt{\textasciitilde{}/MillenniumAnkh/Imscribing/Paraconsistent/}. All theorems are machine-verified, with zero \texttt{sorry} axioms. The project builds with \texttt{lake build Imscribing.Paraconsistent}.

\subsection{Appendix A: Kernel Tuple and Structural Distances}

\begin{table}[htbp]
\centering
\small
\rowcolors{1}{white}{white}
\begin{tabularx}{\textwidth}{>{\centering\arraybackslash}X >{\centering\arraybackslash}X >{\centering\arraybackslash}X}
\toprule
\rowcolor{gray!25}\textbf{Primitive} & \textbf{Value} & \textbf{Description} \\
\midrule
\rowcolors{1}{white}{gray!10}
$\text{Ð}$ (Dimensionality) & $\text{Ð}_{\omega}$ & Holographic, self-written state-space \\
$\text{Þ}$ (Topology) & $\text{Þ}_{\text{O}}$ & Self-referential topology \\
$\text{Ř}$ (Relational) & $\text{Ř}_{\text{=}}$ & Bidirectional feedback \\
$\text{Ф}$ (Parity) & $\text{Ф}_{\text{}}$ & Frobenius-special ($\mu \circ \delta = \text{id}$ exactly) \\
$\text{ƒ}$ (Fidelity) & $\text{ƒ}_{\text{ż}}$ & Quantum coherence essential \\
$\text{Ç}$ (Kinetics) & $\text{Ç}_{\text{@}}$ & Near-equilibrium (slower than observation) \\
$\Gamma$ (Scope) & $\Gamma_{\text{ʔ}}$ & Maximal (aleph) \\
$\text{ɢ}$ (Grammar) & $\text{ɢ}_{\text{ˌ}}$ & Sequential necessity \\
$\text{\ensuremath{\varphî}}$ (Criticality) & $\text{⊙}_{\text{ÿ}}$ & Self-modeling gate open \\
$\text{Ħ}$ (Chirality) & $\text{Ħ}_{\text{A}}$ & 2-step Markov (encounter visible) \\
$\Sigma$ (Stoichiometry) & $\Sigma_{\text{ő}}$ & Many identical components \\
$\Omega$ (Winding) & $\Omega_{\text{z}}$ & Integer winding (topologically protected) \\
\bottomrule
\end{tabularx}
\end{table}

\textbf{Tier:} $\text{O}_{\text{inf}}$ (Frobenius-special) \\
\textbf{Distance to IUG:} 1.3416 (diagonal), 1.7152 (Mahalanobis) \\
\textbf{Differing primitives:} $\Sigma$ ($\Sigma_{\text{ő}}$ vs $\Sigma_{\text{ï}}$), $\text{Ħ}$ ($\text{Ħ}_{\text{A}}$ vs $\text{Ħ}_{\text{!}}$)

\subsection{Appendix B: Theorem Index}

\begin{table}[htbp]
\centering
\small
\rowcolors{1}{white}{white}
\begin{tabularx}{\textwidth}{>{\centering\arraybackslash}X >{\centering\arraybackslash}X >{\centering\arraybackslash}X}
\toprule
\rowcolor{gray!25}\textbf{Theorem} & \textbf{Module} & \textbf{Proof Method} \\
\midrule
\rowcolors{1}{white}{gray!10}
\texttt{no\_explosion} & Belnap & \texttt{simp} \\
\texttt{B\_fixed\_point\_negation} & Belnap & \texttt{rfl} \\
\texttt{B\_is\_top} & Belnap & case analysis \\
\texttt{only\_B\_is\_dialetheic} & DialetheicAlignment & case analysis \\
\texttt{frobenius\_invariant} & Kernel & case analysis on 4 values \\
\texttt{step\_at\_B3} & Kernel & \texttt{simp} \\
\texttt{run\_B3} & Kernel & induction on n \\
\texttt{run\_paradox} & Kernel & induction on n \\
\texttt{complete\_self\_verification} & SelfVerification & composition of prior theorems \\
\texttt{kernel\_is\_O\_inf} & Kernel & \texttt{rfl} \\
\texttt{dialetheic\_alignment} & DialetheicAlignment & composition of prior theorems \\
\texttt{B\_is\_the\_only\_bifurcation\_point} & DialetheicAlignment & \texttt{decide} \\
\texttt{B\_satisfies\_SIC\_axioms} & DialetheicAlignment & composition \\
\texttt{sustain\_preserves\_B} & QCI & induction using \texttt{run\_B3} \\
\texttt{no\_classical\_reduction} & DialetheicAlignment & case analysis \\
\texttt{init\_immortal} & Init & \texttt{Or.inl} \\
\texttt{portal\_type\_is\_O\_inf} & Portal & \texttt{native\_decide} \\
\texttt{paradox\_fs\_is\_O\_inf} & ParadoxFS & \texttt{native\_decide} \\
\texttt{shell\_type\_is\_O\_inf} & ParaconsistentShell & \texttt{native\_decide} \\
\texttt{scheduler\_type\_is\_O\_inf} & CrystalScheduler & \texttt{native\_decide} \\
\texttt{rh\_bridge\_is\_O\_inf} & QCI\_RH\_Bridge & \texttt{decide} \\
\bottomrule
\end{tabularx}
\end{table}

All theorems type-check in Lean 4 with Mathlib v4.28.0. No \texttt{sorry} axioms remain.


\end{document}
